% ch17_adjunctions_deeper.tex — Volume II Chapter 5

\chapter{Adjunctions: A Deeper Account}

\begin{overviewbox}
Chapter~7 introduced adjunctions through examples and established the triangle
identities as a convenient alternative to the hom-set bijection. In this chapter
we take adjunctions apart more carefully. The triangle identities are not merely
\emph{convenient} --- they are \emph{equivalent} to the hom-set bijection, and
that equivalence has a genuine proof. We also show that the right adjoint of a
functor is unique up to natural isomorphism (not merely asserted so), construct
adjunctions from initial objects in comma categories, and ask what happens when
the right adjoint itself has a right adjoint.

By the end you will understand adjunctions from multiple angles
simultaneously: as bijections, as universal morphisms, and as structure in
comma categories. That triangulation is the sign of deep understanding.
\end{overviewbox}


% ═══════════════════════════════════════════════════════════════════════════
\section{Triangle Identities Revisited}
\label{sec:triangle-revisited}
% ═══════════════════════════════════════════════════════════════════════════

\subsection{Informally: What Needs Proof}

Chapter~7 stated that from the natural bijection
$\varphi_{A,B} : \mathcal{D}(FA, B) \xrightarrow{\cong} \mathcal{C}(A, GB)$
one derives a unit $\eta : \mathrm{Id} \Rightarrow GF$ and a counit
$\varepsilon : FG \Rightarrow \mathrm{Id}$ satisfying
\begin{align*}
  (G\varepsilon)(\eta G) &= \mathrm{id}_G, \\
  (\varepsilon F)(F\eta) &= \mathrm{id}_F.
\end{align*}
But we also claimed the converse: given $\eta$ and $\varepsilon$ satisfying
these equations, one can reconstruct the bijection $\varphi$. The full proof
of this equivalence is the business of this section.

\subsection{Expanding: The Two Directions}

The key insight is that $\varphi$ and its inverse can be written entirely in
terms of $\eta$, $\varepsilon$, and the functors:
\begin{align*}
  \varphi(f) &= G(f) \circ \eta_A, \qquad f : FA \to B, \\
  \varphi^{-1}(g) &= \varepsilon_B \circ F(g), \qquad g : A \to GB.
\end{align*}
That these formulas actually define a bijection, and that the bijection is
natural, follow from the triangle identities. Let us verify this in a
structured proof.

\begin{theorem}
\label{thm:triangle-bijection}
Let $F : \mathcal{C} \to \mathcal{D}$ and $G : \mathcal{D} \to \mathcal{C}$
be functors. The following data are equivalent:
\begin{enumerate}
  \item A natural bijection $\varphi_{A,B} : \mathcal{D}(FA, B) \cong \mathcal{C}(A, GB)$.
  \item Natural transformations $\eta : \mathrm{Id}_\mathcal{C} \Rightarrow GF$
        and $\varepsilon : FG \Rightarrow \mathrm{Id}_\mathcal{D}$ satisfying
        the triangle identities.
\end{enumerate}
\end{theorem}

\begin{proof}
We prove that (1) implies (2), and that (2) implies (1), and that the
constructions are mutually inverse.

\textbf{(1) $\Rightarrow$ (2).} Define $\eta_A = \varphi_{A,FA}(\mathrm{id}_{FA})$
and $\varepsilon_B = \varphi^{-1}_{GB,B}(\mathrm{id}_{GB})$. Naturality of
$\varphi$ implies naturality of $\eta$ and $\varepsilon$. The triangle
identities follow from the naturality squares for $\varphi$ evaluated at the
identity morphisms.

\textbf{(2) $\Rightarrow$ (1).} Define $\varphi(f) = G(f) \circ \eta_A$ and
$\psi(g) = \varepsilon_B \circ F(g)$. The triangle identities ensure $\psi
\circ \varphi = \mathrm{id}$ and $\varphi \circ \psi = \mathrm{id}$, as we
verify in the proof log below.
\end{proof}

\subsection{Formalizing: Triangle Identities $\Leftrightarrow$ Bijection}

\begin{prooflog}
\proofstep{Goal: $\psi(\varphi(f)) = f$ for $f : FA \to B$.}{We want to show the
composite is the identity on each morphism.}
\proofstep{$\psi(\varphi(f)) = \varepsilon_B \circ F(\varphi(f))$}{By definition
of $\psi$.}
\proofstep{$= \varepsilon_B \circ F(G(f) \circ \eta_A)$}{Substituting $\varphi(f)
= G(f) \circ \eta_A$.}
\proofstep{$= \varepsilon_B \circ FG(f) \circ F(\eta_A)$}{$F$ is a functor, so
it preserves composition.}
\proofstep{$= f \circ \varepsilon_{FA} \circ F(\eta_A)$}{Naturality of
$\varepsilon$: $\varepsilon_B \circ FG(f) = f \circ \varepsilon_{FA}$.}
\proofstep{$= f \circ \mathrm{id}_{FA} = f$}{The second triangle identity:
$(\varepsilon F)(F\eta) = \mathrm{id}_F$, so $\varepsilon_{FA} \circ F(\eta_A) =
\mathrm{id}_{FA}$.}
\proofstep{Goal: $\varphi(\psi(g)) = g$ for $g : A \to GB$.}{The dual
computation.}
\proofstep{$\varphi(\psi(g)) = G(\psi(g)) \circ \eta_A$}{By definition of
$\varphi$.}
\proofstep{$= G(\varepsilon_B \circ F(g)) \circ \eta_A$}{Substituting $\psi(g)
= \varepsilon_B \circ F(g)$.}
\proofstep{$= G(\varepsilon_B) \circ GF(g) \circ \eta_A$}{$G$ preserves composition.}
\proofstep{$= G(\varepsilon_B) \circ \eta_{GB} \circ g$}{Naturality of $\eta$:
$GF(g) \circ \eta_A = \eta_{GB} \circ g$.}
\proofstep{$= \mathrm{id}_{GB} \circ g = g$}{First triangle identity:
$(G\varepsilon)(\eta G) = \mathrm{id}_G$, so $G(\varepsilon_B) \circ \eta_{GB} =
\mathrm{id}_{GB}$.}
\proofstep{Naturality of $\varphi$.}{For $h : A' \to A$ in $\mathcal{C}$ and
$k : B \to B'$ in $\mathcal{D}$, the naturality squares for $\varphi$ in both
variables follow from the fact that $\eta$ and $\varepsilon$ are natural
transformations and from functoriality of $F$ and $G$.}
\end{prooflog}

\begin{dialog}
Why do we need naturality of $\varepsilon$ in the fourth step above?

The naturality square for $\varepsilon$ at $f : FA \to B$ says
$\varepsilon_B \circ FG(f) = f \circ \varepsilon_{FA}$.
This is exactly the equation that moves $\varepsilon$ past $FG(f)$, trading
the outer $B$-component for the inner $FA$-component. Without it, the proof
would be stuck.
\end{dialog}

\subsection{Reorganizing: Why Both Directions Matter}

The theorem tells us that neither formulation is more fundamental than the
other. In practice:
\begin{itemize}
  \item The \emph{hom-set bijection} form is easiest for computing and
        identifying adjunctions.
  \item The \emph{unit/counit} form is easiest for working in 2-categories
        and proving things about composites of adjunctions.
  \item The formulas $\varphi(f) = Gf \circ \eta$ and $\varphi^{-1}(g) =
        \varepsilon \circ Fg$ let you move between the two whenever needed.
\end{itemize}

\begin{howtosay}
\begin{itemize}
  \item The equations $(G\varepsilon)(\eta G) = \mathrm{id}_G$ and
        $(\varepsilon F)(F\eta) = \mathrm{id}_F$ are called the
        \emph{triangle identities} because they correspond to two triangular
        diagrams (one for each functor).
  \item In the literature you will also see them written as
        $G\varepsilon_B \circ \eta_{GB} = \mathrm{id}_{GB}$ (the $G$-triangle)
        and $\varepsilon_{FA} \circ F\eta_A = \mathrm{id}_{FA}$ (the
        $F$-triangle).
  \item Some authors call these the ``zig-zag identities,'' because the
        composites zig from $\mathcal{C}$ to $\mathcal{D}$ and zag back.
\end{itemize}
\end{howtosay}


% ═══════════════════════════════════════════════════════════════════════════
\section{Uniqueness of Adjoints}
\label{sec:uniqueness-adjoints}
% ═══════════════════════════════════════════════════════════════════════════

\subsection{Informally: Two Right Adjoints Must Be Isomorphic}

Chapter~7 stated without full proof that adjoints are unique up to natural
isomorphism. Now we can prove it rigorously. If $F \dashv G$ and $F \dashv G'$,
then $G \cong G'$ via a canonical natural isomorphism. The proof is the cleanest
illustration of the ``Yoneda argument'': two representable functors agreeing on
all hom-sets must be isomorphic.

\subsection{Expanding: The Idea}

Suppose $G$ and $G'$ are both right adjoints to $F$. For each object $B \in
\mathcal{D}$, we have natural bijections
\begin{align*}
  \mathcal{C}(A, GB) &\cong \mathcal{D}(FA, B), \\
  \mathcal{C}(A, G'B) &\cong \mathcal{D}(FA, B).
\end{align*}
Composing, $\mathcal{C}(A, GB) \cong \mathcal{C}(A, G'B)$ naturally in $A$.
By the Yoneda lemma, a natural isomorphism of representable functors
$\mathcal{C}(-, GB) \cong \mathcal{C}(-, G'B)$ comes from a unique isomorphism
$GB \cong G'B$. Since this works for every $B$, and the isomorphisms assemble
naturally in $B$, we get a natural isomorphism $G \cong G'$.

\begin{theorem}
\label{thm:uniqueness-of-adjoints}
If $F \dashv G$ and $F \dashv G'$, then there is a natural isomorphism
$G \cong G'$. Dually, if $G \dashv F$ and $G' \dashv F$, then $G \cong G'$.
\end{theorem}

\begin{proof}
We construct the isomorphism $\theta : G \Rightarrow G'$ explicitly.
\end{proof}

\begin{prooflog}
\proofstep{For each $B \in \mathcal{D}$, apply both bijections at
$A = G'B$.}{We want a specific isomorphism $GB \to G'B$.}
\proofstep{The bijection for $G'$: $\mathcal{D}(F(G'B), B) \cong
\mathcal{C}(G'B, G'B)$.}{Set $A = G'B$.}
\proofstep{The counit $\varepsilon'_B : F(G'B) \to B$ corresponds to
$\mathrm{id}_{G'B}$ under this bijection.}{Definition of counit as image of
the identity.}
\proofstep{Apply the bijection for $G$ with $A = G'B$, $B = B$:
$\mathcal{D}(F(G'B), B) \cong \mathcal{C}(G'B, GB)$.}{Use the $F \dashv G$
bijection on $\varepsilon'_B$.}
\proofstep{Define $\theta_B := \varphi_G(\varepsilon'_B) = G(\varepsilon'_B)
\circ \eta_{G'B} : G'B \to GB$.}{This is a specific morphism
$G'B \to GB$ for each $B$.}
\proofstep{Symmetrically define $\theta^{-1}_B = G'(\varepsilon_B) \circ
\eta'_{GB} : GB \to G'B$.}{Using the same construction with $G$ and $G'$
swapped.}
\proofstep{$\theta^{-1}_B \circ \theta_B = \mathrm{id}_{G'B}$: compute
$G'(\varepsilon_B) \circ \eta'_{GB} \circ G(\varepsilon'_B) \circ \eta_{G'B}$.}{
Composing the two morphisms.}
\proofstep{By naturality of $\eta'$ and the triangle identity for $G'$
this collapses to $\mathrm{id}_{G'B}$.}{Each triangle identity cancels one
of the zigs.}
\proofstep{Similarly $\theta_B \circ \theta^{-1}_B = \mathrm{id}_{GB}$.}{Symmetric
argument.}
\proofstep{Naturality: for $k : B \to B'$, both sides of $\theta_{B'} \circ
G'(k) = G(k) \circ \theta_B$ follow from naturality of $\eta$, $\varepsilon$,
$\eta'$, $\varepsilon'$.}{Standard naturality check.}
\proofstep{Therefore $\theta : G \Rightarrow G'$ is a natural isomorphism.}{All
components are isomorphisms and assembly is natural.}
\end{prooflog}

\subsection{Reorganizing: The Yoneda Shortcut}

The proof above constructs $\theta_B$ by hand, but the Yoneda lemma offers a
cleaner perspective. The composite
\begin{equation*}
  \mathcal{C}(-, GB) \cong \mathcal{D}(F(-), B) \cong \mathcal{C}(-, G'B)
\end{equation*}
is a natural isomorphism of representable presheaves on $\mathcal{C}$. The
Yoneda lemma says that any natural isomorphism $\mathcal{C}(-, X) \cong
\mathcal{C}(-, Y)$ arises from a unique isomorphism $X \cong Y$. Applied here,
this gives $\theta_B : GB \cong G'B$ for each $B$, naturally in $B$.

The two approaches are equivalent, but the Yoneda version is shorter once you
internalize the lemma.

\subsection{Simplifying: The Slogan}

\begin{equation*}
  \boxed{\;F \dashv G \;\text{ and }\; F \dashv G'
  \;\Longrightarrow\;
  G \cong G' \;\text{(naturally)}.\;}
\end{equation*}

This means we may speak of \emph{the} right adjoint to $F$ without
ambiguity, knowing it is determined up to unique natural isomorphism.


% ═══════════════════════════════════════════════════════════════════════════
\begin{nameorigin}
\term{Comma category} (notation $F \downarrow G$ or $(F, G)$) was
introduced by \emph{F.\ William Lawvere} in his 1963 PhD thesis. The
name reflects the original notation: Lawvere wrote the objects of the
category as triples $(c, d, f)$ separated by commas. The construction
unifies many ad-hoc categories of ``maps with extra data'' (slice
categories, arrow categories, fibrations) under a single universal
construction. Modern category theory uses $\downarrow$ or $/$ notation
to indicate slice-like comma constructions.
\end{nameorigin}

\section{Adjoints as Initial Objects in Comma Categories}
\label{sec:adjoints-initial}
% ═══════════════════════════════════════════════════════════════════════════

\subsection{Informally: The Comma Category Perspective}

There is a third, very geometric way to understand adjunctions. The left
adjoint $F(c)$ is characterized as the \emph{initial object in the comma
category $(c \downarrow G)$}. Let us unpack this.

The comma category $(c \downarrow G)$ has:
\begin{itemize}
  \item Objects: pairs $(d, f)$ where $d \in \mathcal{D}$ and $f : c \to G(d)$
        is a morphism in $\mathcal{C}$.
  \item Morphisms from $(d, f)$ to $(d', f')$: a morphism $h : d \to d'$ in
        $\mathcal{D}$ such that $G(h) \circ f = f'$.
\end{itemize}
An initial object in $(c \downarrow G)$ is an object $(d_0, f_0)$ such that for
every $(d, f)$ there is a unique morphism $(d_0, f_0) \to (d, f)$.

\begin{dialog}
What does an initial object in $(c \downarrow G)$ look like in the free-monoid
example?

Take $c = \{a, b\}$ (a two-element set), $G = $ forgetful functor from monoids
to sets, $\mathcal{D} = \mathbf{Mon}$. The comma category $(c \downarrow G)$ has
objects: pairs $(\text{monoid } M, \text{function } c \to M)$. The initial
object is the free monoid on $\{a, b\}$: it has a canonical function $c \to F(c)$
(send each generator to itself), and any function $c \to M$ factors uniquely
through $F(c)$ by the universal property.
\end{dialog}

\subsection{Expanding: The Theorem}

\begin{theorem}
\label{thm:adjoint-initial}
Let $F : \mathcal{C} \to \mathcal{D}$ and $G : \mathcal{D} \to \mathcal{C}$.
The following are equivalent:
\begin{enumerate}
  \item $F$ is left adjoint to $G$.
  \item For every $c \in \mathcal{C}$, there exists an initial object in the
        comma category $(c \downarrow G)$.
\end{enumerate}
When these hold, the initial object of $(c \downarrow G)$ is $(F(c), \eta_c)$,
where $\eta_c : c \to GF(c)$ is the unit of the adjunction.
\end{theorem}

\begin{proof}
(1) $\Rightarrow$ (2): Given $F \dashv G$, we claim $(F(c), \eta_c)$ is initial
in $(c \downarrow G)$.

Let $(d, f : c \to Gd)$ be any object of $(c \downarrow G)$. By the adjunction
bijection, $f$ corresponds uniquely to a morphism $\tilde{f} : F(c) \to d$ in
$\mathcal{D}$ with $\varphi(\tilde{f}) = f$. That is, $G(\tilde{f}) \circ
\eta_c = f$. This says $\tilde{f}$ is a morphism $(F(c), \eta_c) \to (d, f)$
in $(c \downarrow G)$. Uniqueness of $\tilde{f}$ from bijectivity of $\varphi$
says it is the unique such morphism.

(2) $\Rightarrow$ (1): Given initial objects $(F(c), \eta_c)$ in each $(c
\downarrow G)$, define $\varphi_{c,d}(h) = G(h) \circ \eta_c$ for $h : F(c)
\to d$. Initiality gives a unique $h$ for each $f : c \to Gd$, making $\varphi$
a bijection. Naturality follows from functoriality of $G$ and naturality of $\eta$.

The construction of $F$ on morphisms: for $k : c \to c'$ in $\mathcal{C}$, the
morphism $F(k) : F(c) \to F(c')$ is the unique morphism in $\mathcal{D}$
corresponding (under the adjunction at $c$) to $\eta_{c'} \circ k : c \to GF(c')$.
\end{proof}

\subsection{Formalizing: The Universal Arrow}

The pair $(F(c), \eta_c)$ is called a \term{universal arrow from $c$ to $G$}.
The universal property says: for every $d \in \mathcal{D}$ and morphism $f :
c \to Gd$, there is a unique $\tilde{f} : F(c) \to d$ with $G(\tilde{f}) \circ
\eta_c = f$.

\begin{equation*}
\begin{tikzcd}
  c \arrow[r, "\eta_c"] \arrow[dr, "f"'] & GF(c) \arrow[d, "G(\tilde{f})"] \\
  & Gd
\end{tikzcd}
\quad \leftrightarrow \quad
\begin{tikzcd}
  F(c) \arrow[r, dashed, "\exists!\, \tilde{f}"] & d
\end{tikzcd}
\end{equation*}

This is the ``initial universal arrow'' form of the adjunction. Compare with
the construction of free objects in algebra: the free group on $S$ has a
universal property of exactly this form, with $G$ the forgetful functor.

\subsection{Reorganizing: Three Forms Compared}

\begin{center}
\begin{tabular}{lll}
\textbf{Form} & \textbf{Data} & \textbf{Best for} \\
\hline
Hom-set bijection & $\varphi_{A,B}$ natural & Computing examples \\
Unit/counit & $\eta, \varepsilon$ + triangle ids & 2-categorical arguments \\
Initial in $(c \downarrow G)$ & Universal arrows & Proving existence \\
\end{tabular}
\end{center}

The Adjoint Functor Theorems (Chapter~10) use the third form: to construct a
left adjoint to $G$, show that each comma category $(c \downarrow G)$ has an
initial object.


% ═══════════════════════════════════════════════════════════════════════════
\section{Adjoint of an Adjoint}
\label{sec:adjoint-of-adjoint}
% ═══════════════════════════════════════════════════════════════════════════

\subsection{Informally: Can We Keep Going?}

If $F \dashv G$, can $G$ have a right adjoint too? Certainly. The question is
whether it is related to $F$ in any systematic way, and what the structure of
an adjoint triple looks like.

An \term{adjoint triple} is a sequence $F \dashv G \dashv H$ meaning $F \dashv G$
and $G \dashv H$. In this situation $F$ is a left adjoint of $G$ and $H$ is a
right adjoint of $G$ --- so $G$ is the middle piece with an adjoint on each side.

\subsection{Expanding: Adjoint Triples in Nature}

\begin{example}{Constant Diagram and Limit/Colimit}
Let $\mathcal{J}$ be a small category and consider the functor categories. The
\term{diagonal functor} $\Delta : \mathcal{C} \to \mathcal{C}^\mathcal{J}$ sends
an object $c$ to the constant diagram at $c$. We have:
\begin{equation*}
  \operatorname{colim} \dashv \Delta \dashv \lim,
\end{equation*}
an adjoint triple. The colimit functor is left adjoint to $\Delta$ (since
$\operatorname{colim} D \to c$ corresponds to $D \to \Delta c$), and $\Delta$ is
left adjoint to the limit functor.
\end{example}

\begin{example}{Restriction and Extensions of Sheaves}
For an open inclusion $j : U \hookrightarrow X$ of topological spaces, the
restriction functor $j^* : \mathbf{Sh}(X) \to \mathbf{Sh}(U)$ has both a left
adjoint $j_!$ (extension by zero) and a right adjoint $j_*$ (direct image):
\begin{equation*}
  j_! \dashv j^* \dashv j_*.
\end{equation*}
\end{example}

\begin{example}{Kan Extensions Along a Functor}
For $p : \mathcal{C} \to \mathcal{E}$, the restriction functor
$(-) \circ p : [\mathcal{E}, \mathcal{D}] \to [\mathcal{C}, \mathcal{D}]$ (which
precomposes with $p$) sits in an adjoint triple
\begin{equation*}
  \operatorname{Lan}_p \dashv (-) \circ p \dashv \operatorname{Ran}_p.
\end{equation*}
The left Kan extension is the left adjoint, the right Kan extension is the
right adjoint. This is the content of Chapter~18.
\end{example}

\subsection{Formalizing: Structure of an Adjoint Triple}

\begin{definition}
An \term{adjoint triple} $(F, G, H)$ between categories $\mathcal{C}$ and
$\mathcal{D}$ consists of functors $F, H : \mathcal{C} \to \mathcal{D}$ and
$G : \mathcal{D} \to \mathcal{C}$ (or the transposed version) with $F \dashv G$
and $G \dashv H$.
\end{definition}

An important structural consequence:

\begin{proposition}
In an adjoint triple $F \dashv G \dashv H$, the functor $F$ preserves colimits,
$G$ preserves both limits and colimits, and $H$ preserves limits.
\end{proposition}

\begin{proof}
$F$ is a left adjoint, so it preserves colimits (LAPC). $H$ is a right adjoint,
so it preserves limits (RAPL). $G$ is simultaneously a right adjoint (to $F$,
so RAPL applies) and a left adjoint (to $H$, so LAPC applies). Thus $G$ preserves
both.
\end{proof}

\subsection{Reorganizing: The Double Right Adjoint}

Starting from $F \dashv G$, the ``double right adjoint'' question asks: is
$G \dashv ?$. The answer depends entirely on the specific functors involved ---
there is no general construction of a right adjoint to $G$ from $F$ alone. The
existence of a right adjoint to $G$ is an independent piece of structure.

What \emph{can} be said: by uniqueness of adjoints (Theorem~\ref{thm:uniqueness-of-adjoints}),
if $G$ has any right adjoint, that right adjoint is unique up to natural
isomorphism. So the set of adjoint triples containing a given $F \dashv G$ is
either empty or essentially unique.

\subsection{Simplifying: Higher Adjoints}

One can ask about adjoint \emph{quadruples}, quintuples, and so on. In good
situations (locally presentable categories, Grothendieck toposes) these finite
adjoint towers arise naturally. In the theory of $\infty$-categories there are
even infinite adjoint chains. The pattern of limit/colimit preservation escalates
accordingly.

\begin{dialog}
Is there a way to tell in advance whether $G$ has a right adjoint, without
constructing one?

Yes: the Adjoint Functor Theorems (Chapter~10) give conditions. Roughly, $G$
has a right adjoint iff it preserves colimits and a smallness/solution-set
condition holds. The theorem tells you to check preservation rather than construct.
\end{dialog}


% ═══════════════════════════════════════════════════════════════════════════
\section*{Exercises}
\addcontentsline{toc}{section}{Exercises}
% ═══════════════════════════════════════════════════════════════════════════

\begin{exercisesbox}
\begin{qa}
\qaitem{State the triangle identities in terms of the unit and counit.}{$(G\varepsilon)(\eta G) = \mathrm{id}_G$ and $(\varepsilon F)(F\eta) = \mathrm{id}_F$.}
\qaitem{Write the adjunction bijection $\varphi$ in terms of unit and counit.}{$\varphi(f) = G(f) \circ \eta_A$ for $f : FA \to B$.}
\qaitem{Write $\varphi^{-1}$ in terms of unit and counit.}{$\varphi^{-1}(g) = \varepsilon_B \circ F(g)$ for $g : A \to GB$.}
\qaitem{Which naturality of $\varepsilon$ is used in the proof that $\psi \circ \varphi = \mathrm{id}$?}{The naturality $\varepsilon_B \circ FG(f) = f \circ \varepsilon_{FA}$.}
\qaitem{What does it mean for $\psi \circ \varphi = \mathrm{id}$?}{The formula $\varphi^{-1}(\varphi(f)) = f$ for all $f : FA \to B$.}
\qaitem{Suppose $F \dashv G$ and $F \dashv G'$. What is the canonical isomorphism $\theta_B : G'B \to GB$?}{$\theta_B = G(\varepsilon'_B) \circ \eta_{G'B}$, using the $G'$-counit applied via the $G$-unit.}
\qaitem{What is the Yoneda argument for uniqueness of adjoints?}{The bijection $\mathcal{C}(-,GB) \cong \mathcal{D}(F-,B) \cong \mathcal{C}(-,G'B)$ is a natural iso of representables, so the Yoneda lemma gives $GB \cong G'B$.}
\qaitem{What is the comma category $(c \downarrow G)$?}{Objects are pairs $(d, f : c \to Gd)$; morphisms are $h : d \to d'$ with $G(h) \circ f = f'$.}
\qaitem{What is the initial object in $(c \downarrow G)$ when $F \dashv G$ exists?}{$(F(c), \eta_c : c \to GF(c))$.}
\qaitem{What is a universal arrow from $c$ to $G$?}{A pair $(F(c), \eta_c)$ such that every $f : c \to Gd$ factors uniquely as $G(\tilde{f}) \circ \eta_c$.}
\qaitem{Reformulate the adjunction $F \dashv G$ purely using universal arrows.}{For every $c \in \mathcal{C}$, there is a universal arrow from $c$ to $G$.}
\qaitem{Give an example of an adjoint triple involving the diagonal functor.}{$\operatorname{colim} \dashv \Delta \dashv \lim$.}
\qaitem{In $j_! \dashv j^* \dashv j_*$ for an open inclusion $j$, which functor preserves both limits and colimits?}{The restriction $j^*$, being both a left and right adjoint.}
\qaitem{In an adjoint triple $F \dashv G \dashv H$, which functor preserves limits?}{Both $G$ (as right adjoint to $F$) and $H$ (as right adjoint to $G$).}
\qaitem{Does $F \dashv G$ force $G$ to have a right adjoint?}{No; the existence of a right adjoint to $G$ is independent structure.}
\qaitem{If $G$ has a right adjoint, how many essentially different ones does it have?}{Exactly one, up to natural isomorphism (by uniqueness of adjoints).}
\qaitem{What does RAPL say about the middle functor in an adjoint triple?}{It preserves limits (being a right adjoint) and colimits (being a left adjoint).}
\qaitem{What is the adjoint triple for Kan extensions along $p : \mathcal{C} \to \mathcal{E}$?}{$\operatorname{Lan}_p \dashv (-) \circ p \dashv \operatorname{Ran}_p$.}
\qaitem{What is the ``solution-set condition'' used for?}{It is a smallness hypothesis in the Adjoint Functor Theorem that, combined with limit-preservation, guarantees the existence of a left adjoint.}
\qaitem{Are there categories admitting infinite adjoint chains?}{Yes: in sufficiently well-behaved categories (like presheaf toposes) one finds arbitrarily long adjoint chains.}
\end{qa}
\end{exercisesbox}


% ═══════════════════════════════════════════════════════════════════════════
\section*{Chapter Summary}
\addcontentsline{toc}{section}{Chapter Summary}
% ═══════════════════════════════════════════════════════════════════════════

\begin{summarybox}
\textbf{Triangle identities $\Leftrightarrow$ hom-set bijection.} The formulas
$\varphi(f) = G(f) \circ \eta_A$ and $\varphi^{-1}(g) = \varepsilon_B \circ F(g)$
convert between the two presentations of an adjunction. The triangle identities
are exactly the condition that these formulas define inverse bijections.

\textbf{Uniqueness.} If $F \dashv G$ and $F \dashv G'$, then $G \cong G'$
canonically. The canonical isomorphism $\theta_B$ is constructed by applying
one adjunction bijection to the counit of the other.

\textbf{Comma categories.} The left adjoint $F(c)$ is the initial object in
the comma category $(c \downarrow G)$. Equivalently, $(F(c), \eta_c)$ is the
universal arrow from $c$ to $G$. This third formulation of adjunctions is the
basis for the Adjoint Functor Theorems.

\textbf{Adjoint triples.} When $G$ itself has a right adjoint $H$, we get
$F \dashv G \dashv H$. The middle functor $G$ then preserves both limits and
colimits. Examples: $\operatorname{colim} \dashv \Delta \dashv \lim$, and Kan
extensions $\operatorname{Lan}_p \dashv (-) \circ p \dashv \operatorname{Ran}_p$.
\end{summarybox}


% ═══════════════════════════════════════════════════════════════════════════
\section*{Formal Restatement}
\addcontentsline{toc}{section}{Formal Restatement}
% ═══════════════════════════════════════════════════════════════════════════

\begin{formalrestatement}
\textbf{Triangle identities.} Given $F : \mathcal{C} \to \mathcal{D}$,
$G : \mathcal{D} \to \mathcal{C}$, $\eta : \mathrm{Id}_\mathcal{C} \Rightarrow GF$,
$\varepsilon : FG \Rightarrow \mathrm{Id}_\mathcal{D}$, the following are
equivalent: (i) $F \dashv G$ with unit $\eta$ and counit $\varepsilon$; (ii)
$(G\varepsilon)(\eta G) = \mathrm{id}_G$ and $(\varepsilon F)(F\eta) =
\mathrm{id}_F$. Under these conditions, $\varphi_{A,B}(f) = G(f) \circ \eta_A$
and $\varphi^{-1}_{A,B}(g) = \varepsilon_B \circ F(g)$ are mutually inverse
natural bijections.

\medskip
\textbf{Uniqueness.} If $F \dashv G$ and $F \dashv G'$, then $G \cong G'$
via $\theta_B = G(\varepsilon'_B) \circ \eta_{G'B}$, a natural isomorphism.

\medskip
\textbf{Comma category characterization.} $F \dashv G$ iff for each $c$,
the comma category $(c \downarrow G)$ has an initial object, which is
$(F(c), \eta_c : c \to GF(c))$.

\medskip
\textbf{Adjoint triple.} $F \dashv G \dashv H$ means $F \dashv G$ and
$G \dashv H$. The middle functor $G$ preserves limits (RAPL applied to
$F \dashv G$) and colimits (LAPC applied to $G \dashv H$).

\medskip
\noindent\textit{Chapter~18 develops Kan extensions as adjoint triples and
shows how every adjunction can be expressed using Kan extensions along the
Yoneda embedding.}
\end{formalrestatement}
